\Needspace{7\baselineskip}\Article{298}
\textsc{Aufgabe.} \textit{Es seien \(a,b,c\) beliebige Zahlen, von denen keine \(=0\) ist. Gesucht sind die Bedingungen für die Lösbarkeit der Gleichung}

\[axx+byy+czz=0\ldots(\omega).\]
\textit{Lösung.} Es seien \(\alpha\alpha,\beta\beta,\gamma\gamma\) die größten Quadrate, die der Reihe nach \(bc,ac,ab\) teilen, und es sei \(\alpha a=\beta\gamma A,\ \beta b=\alpha\gamma B,\ \gamma c=\alpha\beta C\). Dann sind \(A,B,C\) paarweise teilerfremde ganze Zahlen. Die Gleichung \((\omega)\) ist genau dann lösbar, wenn
\[AXX+BYY+CZZ=0\ldots(\Omega)\]
lösbar ist. Das lässt sich nach Art.\ 294 entscheiden.

\textit{Beweis.} Man setze \(bc=\mathfrak A\alpha\alpha,\ ac=\mathfrak B\beta\beta,\ ab=\mathfrak C\gamma\gamma\). Dann sind \(\mathfrak A,\mathfrak B,\mathfrak C\) ganze Zahlen ohne quadratische Faktoren, und es gilt \(\mathfrak A=BC,\ \mathfrak B=AC,\ \mathfrak C=AB\). Daraus folgt \(\mathfrak A\mathfrak B\mathfrak C=(ABC)^2\), also ist \(ABC=A\mathfrak A=B\mathfrak B=C\mathfrak C\) notwendig eine ganze Zahl. Der größte gemeinsame Teiler der Zahlen \(\mathfrak A,A\mathfrak A\) sei \(m\), und es sei \(\mathfrak A=gm,\ A\mathfrak A=hm\). Dann ist \(g\) teilerfremd zu \(h\) und, da \(\mathfrak A\) keine quadratischen Faktoren hat, auch zu \(m\). Nun gilt \(hhm=gAA\mathfrak A =g\mathfrak B\mathfrak C\), also teilt \(g\) die Zahl \(hhm\). Das ist offenbar nur möglich, wenn \(g=\pm1\). Somit gilt \(\mathfrak A=\pm m,\ A=\pm h\), also ist diese Zahl ganz; ebenso sind \(B,C\) ganz. \textit{Damit ist die erste Behauptung bewiesen.} — Da \(\mathfrak A=BC\) keine quadratischen Faktoren enthält, müssen \(B,C\) teilerfremd sein. Ebenso ist \(A\) teilerfremd zu \(C\) und zu \(B\). \textit{Damit ist die zweite Behauptung bewiesen.} — Wird schließlich \((\Omega)\) durch \(X=P,Y=Q\), \(Z=R\) erfüllt, so wird offenbar \((\omega)\) durch \(x=\alpha P,y=\beta Q,z=\gamma R\) gelöst. Wird umgekehrt diese durch \(x=p,y=q,z=r\) erfüllt, so wird jene durch \(X=\beta\gamma p\), \(Y=\alpha\gamma q,Z=\alpha\beta r\) erfüllt. Also sind entweder beide lösbar oder keine. \textit{Damit ist die dritte Behauptung bewiesen.}

\Needspace{9\baselineskip}\begin{center}Darstellung der Null durch beliebige ternäre Formen.\end{center}


\Needspace{7\baselineskip}\Article{299}
\textsc{Aufgabe.} \textit{Für eine gegebene ternäre Form}

\[f=axx+a'x'x'+a''x''x''+2bx'x''+2b'xx''+2b''xx'\]

\SourceLeaf{364}
\textit{ist zu entscheiden, ob sie die Null darstellen kann, wobei die Werte der Unbestimmten nicht alle \(=0\) sein dürfen.}

\textit{Lösung.} I.\ Ist \(a=0\), so können die Werte von \(x^\prime,x^{\prime\prime}\) beliebig gewählt werden. Aus der Gleichung
\[a^\prime x^\prime x^\prime+2b x^\prime x^{\prime\prime}+a^{\prime\prime}x^{\prime\prime}x^{\prime\prime}=-2x(b^\prime x^{\prime\prime}+b^{\prime\prime}x^\prime)\]
ergibt sich offenbar ein bestimmter rationaler Wert für \(x\). Ergibt sich auf diese Weise für \(x\) ein Bruch, so braucht man die Werte von \(x,x^\prime,x^{\prime\prime}\) nur mit dessen Nenner zu multiplizieren, um ganze Zahlen zu erhalten. Auszuschließen sind nur die Werte von \(x^\prime,x^{\prime\prime}\), die \(b^\prime x^{\prime\prime}+b^{\prime\prime}x^\prime=0\) ergeben, sofern sie nicht zugleich \(a^\prime x^\prime x^\prime+2b x^\prime x^{\prime\prime}+a^{\prime\prime}x^{\prime\prime}x^{\prime\prime}=0\) ergeben; in diesem Fall kann \(x\) beliebig gewählt werden. Offenbar erhält man so alle möglichen Lösungen. Der Fall, dass \(b^\prime\) null ist und \(b^{\prime\prime}=0\) gilt, gehört allerdings nicht hierher: Dann kommt \(x\) in \(f\) nicht vor, also ist \(f\) eine binäre Form. Ob die Null durch \(f\) darstellbar ist, muss dann nach der Theorie dieser Formen entschieden werden.

II.\ Gilt dagegen nicht \(a=0\), so ist die Gleichung \(f=0\) gleichbedeutend mit
\[(ax+b^{\prime\prime}x^\prime+b^\prime x^{\prime\prime})^2-A^{\prime\prime}x^\prime x^\prime+2B x^\prime x^{\prime\prime}-A^\prime x^{\prime\prime}x^{\prime\prime}=0\]
wobei man
\[\begin{gathered}b^{\prime\prime}b^{\prime\prime}-aa^\prime=A^{\prime\prime},\\ ab-b^\prime b^{\prime\prime}=B,\\ b^\prime b^\prime-aa^{\prime\prime}=A^\prime\end{gathered}\]
setzt. Ist hier \(A^\prime=0\), aber nicht \(B=0\), so ergeben sich bei beliebiger Wahl von \(ax+b^{\prime\prime}x^\prime+b^\prime x^{\prime\prime}\) und \(x^{\prime\prime}\) offenbar rationale Werte für \(x\) und \(x^\prime\). Sind sie nicht ganz, so werden sie wenigstens durch einen geeigneten Multiplikator ganz. Für den einen Wert von \(x^{\prime\prime}\), nämlich \(x^{\prime\prime}=0\), ist der Wert von \(ax+b^{\prime\prime}x^\prime+b^\prime x^{\prime\prime}\) nicht beliebig, sondern muss ebenfalls \(=0\) gesetzt werden. Dann kann aber \(x^\prime\) beliebig gewählt werden und liefert einen rationalen Wert für \(x\). — Ist dagegen zugleich \(A^{\prime\prime}\) null und gilt \(B=0\), und ist \(A^\prime\) ein Quadrat \(=kk\), so zerfällt \(f=0\) offenbar in die beiden linearen Gleichungen
\[\begin{gathered}ax+b^{\prime\prime}x^\prime+(b^\prime+k)x^{\prime\prime}=0,\\ ax+b^{\prime\prime}x^\prime+(b^\prime-k)x^{\prime\prime}=0\end{gathered}\]
von denen \textit{entweder} die eine \textit{oder} die andere gelten muss. Ist hingegen unter derselben Voraussetzung \(A^\prime\) kein Quadrat, so hängt die Lösung der vorgelegten Gleichung offenbar von \(x^{\prime\prime}=0\) und \(ax+b^{\prime\prime}x^\prime=0\) ab, die \textit{zugleich} gelten müssen.
\SourceLeaf{365}

Es braucht kaum erwähnt zu werden, dass die Methode in I auch anwendbar ist, wenn \(a^\prime\) null ist oder \(a^{\prime\prime}=0\) gilt, und die Methode in II, wenn \(A^\prime=0\).

III.\ Ist dagegen weder \(a\) null, noch gilt \(A^{\prime\prime}=0\), so ist \(f=0\) gleichbedeutend mit
\[A^{\prime\prime}(ax+b^{\prime\prime}x^\prime+b^\prime x^{\prime\prime})^2-(A^{\prime\prime}x^\prime-Bx^{\prime\prime})^2+Da x^{\prime\prime}x^{\prime\prime}=0\]
wobei \(D\) die Determinante der Form \(f\) bezeichnet oder \(Da\) die Zahl \(BB-A^\prime A^{\prime\prime}\). Ist \(D=0\), so verhält sich die Lösung ähnlich wie am Ende des vorigen Falls. Ist nämlich \(A^{\prime\prime}\) ein Quadrat \(=kk\), so zerfällt die vorgelegte Gleichung in
\[\begin{gathered}kax+(kb^{\prime\prime}-A^{\prime\prime})x^\prime+(kb^\prime+B)x^{\prime\prime}=0,\\ kax+(kb^{\prime\prime}+A^{\prime\prime})x^\prime+(kb^{\prime\prime}-B)x^{\prime\prime}=0\end{gathered}\]
ist dagegen \(A^{\prime\prime}\) kein Quadrat, so muss
\[\begin{gathered}ax+b^{\prime\prime}x^\prime+b^\prime x^{\prime\prime}=0,\\ A^{\prime\prime}x^\prime-Bx^{\prime\prime}=0\end{gathered}\]
gelten. Ist aber \(D\) nicht \(=0\), so sind wir auf die Gleichung
\[A^{\prime\prime}tt-uu+Davv=0\]
geführt, deren Lösbarkeit sich nach dem vorigen Artikel entscheiden lässt. Hat diese nur die Lösung \(t=0,\ u=0,\ v=0\), so hat offenbar auch die vorgelegte keine andere Lösung als \(x=0,\ x^\prime=0,\ x^{\prime\prime}=0\). Ist jene aber anders lösbar, so ergeben sich aus beliebigen solchen ganzzahligen Werten von \(t,u,v\) durch die Gleichungen
\[\begin{gathered}ax+b^{\prime\prime}x^\prime+b^\prime x^{\prime\prime}=t,\\ A^{\prime\prime}x^\prime-Bx^{\prime\prime}=u,\\ x^{\prime\prime}=v\end{gathered}\]
wenigstens rationale Werte von \(x,x^\prime,x^{\prime\prime}\). Enthalten diese Brüche, so erhält man durch einen geeigneten Multiplikator ganze Zahlen.

Sobald aber \textit{eine} ganzzahlige Lösung von \(f=0\) gefunden ist, kann man die Aufgabe auf Fall I zurückführen und wie dort sämtliche Lösungen auf folgende Weise angeben. Die Gleichung \(f=0\) werde durch die Werte von \(x,x^\prime,x^{\prime\prime}\) erfüllt, nämlich \(\alpha,\alpha^\prime,\alpha^{\prime\prime}\), die wir ohne gemeinsamen Faktor voraussetzen. Nach Art.\ 40, 279 wähle man ganze Zahlen \(\beta,\beta^\prime,\beta^{\prime\prime},\ \gamma,\gamma^\prime,\gamma^{\prime\prime}\) so, dass
\[\alpha(\beta^\prime\gamma^{\prime\prime}-\beta^{\prime\prime}\gamma^\prime)+\alpha^\prime(\beta^{\prime\prime}\gamma-\beta\gamma^{\prime\prime})+\alpha^{\prime\prime}(\beta\gamma^\prime-\beta^\prime\gamma)=1\]

\SourceLeaf{366}
und es gehe \(f\) durch die Substitution
\[\begin{gathered}x=\alpha y+\beta y^\prime+\gamma y^{\prime\prime}.\\ x^\prime=\alpha^\prime y+\beta^\prime y^\prime+\gamma^\prime y^{\prime\prime},\\ x^{\prime\prime}=\alpha^{\prime\prime}y+\beta^{\prime\prime}y^\prime+\gamma^{\prime\prime}y^{\prime\prime}\ldots(S)\end{gathered}\]
in
\[g=cyy+c^\prime y^\prime y^\prime+c^{\prime\prime}y^{\prime\prime}y^{\prime\prime}+2dy^\prime y^{\prime\prime}+2d^\prime yy^{\prime\prime}+2d^{\prime\prime}yy^\prime\]
über. Dann ist offenbar \(c=0\), und \(g\) ist zu \(f\) äquivalent. Daher ergeben sich aus allen Lösungen von \(g=0\) durch \(S\) sämtliche ganzzahligen Lösungen von \(f=0\). Nach I sind nun alle Lösungen von \(g=0\) in den Formeln
\[\begin{gathered}y=-z(c^\prime pp+2dpq+c^{\prime\prime}qq),\\ y^\prime=2z(d^{\prime\prime}pp+d^\prime pq),\\ y^{\prime\prime}=2z(d^{\prime\prime}pq+d^\prime qq)\end{gathered}\]
enthalten. Dabei bezeichnen \(p,q\) beliebige ganze Zahlen und \(z\) eine beliebige Zahl, für die auch Brüche zugelassen sind, sofern \(y,y^\prime,y^{\prime\prime}\) ganz bleiben. Setzt man diese Werte von \(y,y^\prime,y^{\prime\prime}\) in \((S)\) ein, so erhält man alle ganzzahligen Lösungen von \(f=0\). — Ist zum Beispiel
\[f=xx+x^\prime x^\prime+x^{\prime\prime}x^{\prime\prime}-4x^\prime x^{\prime\prime}+2xx^{\prime\prime}+8xx^\prime\]
und hat man eine Lösung von \(f=0\), nämlich \(x=1,\ x^\prime=-2,\ x^{\prime\prime}=1\), so ergibt sich für \(\beta,\beta^\prime,\beta^{\prime\prime},\ \gamma,\gamma^\prime,\gamma^{\prime\prime}=0,1,0,0,0,1\)
\[g=y^\prime y^\prime+y^{\prime\prime}y^{\prime\prime}-4y^\prime y^{\prime\prime}+12yy^{\prime\prime}\]
Damit sind sämtliche ganzzahligen Lösungen von \(g=0\) in der Formel
\[\begin{gathered}y=-z(pp-4pq+qq),\\ y^\prime=12zpq,\\ y^{\prime\prime}=12zqq\end{gathered}\]
enthalten und somit alle Lösungen von \(f=0\) in
\[\begin{gathered}x=-z(pp-4pq+qq)\\x^\prime=2z(pp+2pq+qq)\\x^{\prime\prime}=-z(pp-4pq-11qq)\end{gathered}\]


\Needspace{9\baselineskip}\begin{center}Allgemeine Lösung unbestimmter Gleichungen zweiten Grades mit zwei Unbekannten durch rationale Zahlen.\end{center}


\Needspace{7\baselineskip}\Article{300}
Aus der Aufgabe des vorigen Artikels ergibt sich unmittelbar die Lösung der unbestimmten Gleichung
\[axx+2bxy+cyy+2dx+2ey+f=0\]

\SourceLeaf{367}
wenn nur rationale Werte verlangt werden. Für ganzzahlige Werte haben wir sie schon oben in Art.\ 216 und den folgenden Artikeln gelöst. Denn sämtliche rationalen Werte von \(x,y\) lassen sich als \(\frac{t}{v},\frac{u}{v}\) schreiben, wobei \(t,u,v\) ganz sind. Die Lösung jener Gleichung durch rationale Zahlen ist also offenbar gleichbedeutend mit der ganzzahligen Lösung von
\[att+2btu+cuu+2dtv+2euv+fvv=0\]
und diese stimmt mit der im vorigen Artikel behandelten Gleichung überein. Auszuschließen sind nur die Lösungen mit \(v=0\); solche können aber nicht auftreten, wenn \(bb-ac\) kein Quadrat ist. So sind zum Beispiel alle rationalen Lösungen der Gleichung
\[xx+8xy+yy+2x-4y+1=0\]
die in Art.\ 221 allgemein durch ganze Zahlen gelöst wurde, in der Formel
\[\begin{gathered}x=\frac{pp-4pq+qq}{pp-4pq-11qq},\\ y=-\frac{2pp+4pq+2qq}{pp-4pq-11qq}\end{gathered}\]
enthalten, wobei \(p,q\) beliebige ganze Zahlen sind. — Wir haben diese beiden eng verbundenen Aufgaben hier nur kurz behandelt und viele zugehörige Bemerkungen weggelassen, teils um nicht zu ausführlich zu werden, teils weil wir eine andere Lösung der Aufgabe des vorigen Artikels haben, die auf allgemeineren Grundsätzen beruht. Da ihre Darstellung eine tiefere Untersuchung der ternären Formen verlangt, müssen wir sie einer anderen Gelegenheit vorbehalten.
